% =====================================================================
% Structured summary (at most 3 pages), shared by Oriana's PRDL-DL and
% MLLB-DL reports; the last paragraph depends on \ifprdl.
% Required by the professor's message of 20/11/2026: scope, core methodology,
% main findings, key conclusions.
% =====================================================================
\section*{Structured summary}
\addcontentsline{toc}{section}{Structured summary}

\paragraph{Scope.}
I applied the deep-learning modules of the course to a single clinical problem: deciding which of the
five diagnostic superclasses of PTB-XL (NORM, MI, STTC, CD, HYP) a ten-second 12-lead ECG carries
\cite{wagner2020ptbxl,ptbxl_physionet,pollard2026physionet}. The dataset has \edaNRecordsAll{} records
from \edaNPatientsAll{} patients, used at 100\,Hz, and a record can belong to several superclasses. The
three areas of the guidelines \cite{course_message_2011} are covered with models small enough for a
4\,GB laptop GPU: core architectures (dense network, 1D CNN, LSTM/GRU, RNN with attention,
self-attention and a Transformer encoder); learning strategies (our own label-free ECG-GPT, reused by
linear probing and fine-tuning, and CLIP used zero-shot); and generative AI (an ECG-GPT and an LSTM
that generate lead II sample by sample). We were a team of two. I owned the exploratory analysis, the
dense network, the CNN and the generators; Mar\'ia owned the data pipeline, the evaluation protocol,
the recurrent networks and the Transformer, which I summarise and attribute; the learning-strategy
experiments and the final comparison we did together.

\paragraph{Core methodology.}
\begin{itemize}[leftmargin=1.2em,itemsep=1pt,topsep=2pt]
  \item \textbf{Split and input.} The patient-wise folds of the dataset (1--8 for training, 9 for
        validation, 10 for test) \cite{strodthoff2021deep}; the \edaNNosuperAll{} records with no
        diagnostic superclass are excluded, as the benchmark does \cite{strodthoff2021deep}; each lead is
        z-scored with statistics of the training records.
  \item \textbf{Measuring.} The main score is the macro AUROC (mean of the five per-class AUROCs, no
        threshold needed); macro F1 uses per-class thresholds fixed on validation.
  \item \textbf{Deciding.} Every block was planned in writing before its runs: one factor changed at a
        time, three seeds per configuration, ablations and the choice of the final model on validation
        with a rule set in advance, one test evaluation per final model. Two models are compared on the
        same patient-level bootstrap resamples; I call an effect consistent only when the 95\,\%
        interval of the paired difference excludes 0 in all three seeds.
  \item \textbf{Traceability.} Every number printed here is generated from the run logs, and every
        statement attributed to a paper rests on a quote checked against its PDF. Code and logs:
        \url{https://github.com/orianafe16/ptbxl-DL}.
\end{itemize}

\paragraph{Main findings.} Table~\ref{tab:summary} lists the final models (test macro AUROC, mean
$\pm$ SD over three seeds).

\begin{table}[H]
\centering
\caption{Final models on the test fold (macro AUROC, mean $\pm$ SD over seeds).}
\label{tab:summary}
\small
\begin{tabular}{lcc}
\toprule
Model & 12 leads & Lead II only \\
\midrule
Trivial (logistic regression on lead means and SDs) & \comparisonTestAurocMeanTrivial{} & --- \\
Dense network & \comparisonTestAurocMeanTwelveDense{} $\pm$ \comparisonTestAurocSdTwelveDense{} & \comparisonTestAurocMeanOneleadDense{} $\pm$ \comparisonTestAurocSdOneleadDense{} \\
1D CNN & \textbf{\comparisonTestAurocMeanTwelveCnn{}} $\pm$ \comparisonTestAurocSdTwelveCnn{} & \textbf{\comparisonTestAurocMeanOneleadCnn{}} $\pm$ \comparisonTestAurocSdOneleadCnn{} \\
RNN (conv.\ front end + LSTM) & \comparisonTestAurocMeanTwelveRnn{} $\pm$ \comparisonTestAurocSdTwelveRnn{} & \comparisonTestAurocMeanOneleadRnn{} $\pm$ \comparisonTestAurocSdOneleadRnn{} \\
RNN + attention & \comparisonTestAurocMeanTwelveRnnattn{} $\pm$ \comparisonTestAurocSdTwelveRnnattn{} & \comparisonTestAurocMeanOneleadRnnattn{} $\pm$ \comparisonTestAurocSdOneleadRnnattn{} \\
Transformer encoder & \comparisonTestAurocMeanTwelveTransformer{} $\pm$ \comparisonTestAurocSdTwelveTransformer{} & \comparisonTestAurocMeanOneleadTransformer{} $\pm$ \comparisonTestAurocSdOneleadTransformer{} \\
\midrule
ECG-GPT from scratch, 100\,\% / 10\,\% of labels & --- & \comparisonTestAurocMeanScratchFOneZeroZero{} / \comparisonTestAurocMeanScratchFOneZero{} \\
ECG-GPT linear probe, 100\,\% / 10\,\% & --- & \comparisonTestAurocMeanProbeFOneZeroZero{} / \comparisonTestAurocMeanProbeFOneZero{} \\
ECG-GPT fine-tuning, 100\,\% / 10\,\% & --- & \comparisonTestAurocMeanFinetuneFOneZeroZero{} / \comparisonTestAurocMeanFinetuneFOneZero{} \\
CLIP zero-shot (lead II drawn as an image) & --- & \comparisonTestAurocMeanZeroshot{} \\
\midrule
Published best single model, 12 leads, other setting \cite{strodthoff2021deep} & \comparisonPublishedStrodthoffResnetOneDWang{} & --- \\
\bottomrule
\end{tabular}
\end{table}

\begin{itemize}[leftmargin=1.2em,itemsep=1pt,topsep=2pt]
  \item \textbf{Area 1: weight sharing over time is the step that matters.} The dense network, which has
        one weight per lead and instant, barely beats the trivial model; its result depends mostly on
        depth, batch norm and dropout. The CNN is \cnnDiffDenseMean{} above it in every seed with far fewer
        parameters, and none of its ablations moves validation by more than a few thousandths. Recurrent
        and attention models, which also share weights over time, come close to the CNN but stay below it
        in every seed (RNN $-$ CNN \rnnDiffCnnMean{}, Transformer $-$ CNN \transformerDiffCnnMean{}).
        Mar\'ia's analysis shows that attention over the RNN states behaves like a learned mean and that
        the Transformer ignores the order of its patches.
  \item \textbf{Lead II alone loses information.} Every family drops in every seed when it sees only
        lead II, by \leadsTestDiffMeanTransformer{} (Transformer) to \leadsTestDiffMeanRnnattn{}
        (RNN + attention).
  \item \textbf{Area 2: pre-training pays off only when the network is fine-tuned.} The ECG-GPT, trained
        without labels on the training folds to predict the next sample, gains over the same network
        trained from scratch in every seed; fine-tuned with 10\,\% of the training labels it even beats
        scratch with all of them (\transferTestDiffMeanFinetuneFOneZeroVsScratchFOneZeroZero{}), although
        early stopping and thresholds still used the full validation fold. A $2\times2$ ablation
        attributes the gain to the pre-trained weights, not to the smaller learning rate. Frozen, its
        features are weak, and the generator chosen for generating best gave the weakest ones. On lead II
        a CNN from scratch remains better (\comparisonTestDiffMeanFinetuneFOneZeroZeroVsOneleadCnn{}).
        CLIP works zero-shot on CIFAR-100 (top-1 \zeroshotCifarTopOneNotebook{}, chance
        \zeroshotCifarChanceTopOne{}) and recognises ECG images as ECGs, but its diagnosis is near chance
        (\zeroshotTestMacroAuroc{}), with MI ranked against the truth.
  \item \textbf{Area 3: the generators learn the ECG, and decoding decides what comes out.} With the
        causal mask every generator needs fewer bits per sample than a bigram counting model (baseline
        \generativeBitsMeanBase{} against \generativeBigramBitsVTwoFiveSix{}). The same-size LSTM, a
        longer context and finer tokens do better; without positional encoding the ECG-GPT does worse,
        the reverse of the classifier. Greedy decoding gives a flat line in every record, temperature
        0.8--1 gives beats with a realistic amplitude distribution but an imprecise rate, and
        temperature 1.5 gives noise.
\end{itemize}

\paragraph{Key conclusions.}
(i) On a dataset of this size, the inductive bias of local filters shared over time matters more than
the family name or than self-supervised pre-training. (ii) Inside a sound architecture most ablations
change the result little, and several final models picked by our rules cannot be told apart from their
baselines; we report them as such. (iii) Pre-training without labels makes the network much more
label-efficient, but the best generator was not the best feature extractor. (iv) Zero-shot
classification needs a shared input--text space that covers the task, not only the modality.
(v) A generator must be judged with a vocabulary-independent metric and by its samples; the decoding
rule matters as much as the model. Limits: one dataset and five labels, small models with a 50-epoch
budget (the generators had not converged), a ``foundation'' model pre-trained on PTB-XL only, and no
external or clinical validation.
\ifprdl
\par\smallskip\noindent\textit{This PRDL report} explains the theory behind each module and analyses the
results critically; the code, training details and PyTorch implementation are in the companion
MLLB-DL report.
\else
\par\smallskip\noindent\textit{This MLLB report} describes how the experiments were designed,
implemented in PyTorch, trained, evaluated and made reproducible; the theory and the critical analysis
are in the companion PRDL-DL report.
\fi
